\documentclass[11pt,a4paper]{article}
\usepackage[margin=2cm]{geometry}
\usepackage[T1]{fontenc}
\usepackage{lmodern}
\usepackage{booktabs}
\usepackage{enumitem}
\usepackage[hidelinks]{hyperref}
\setlist{nosep,leftmargin=*}
\setlength{\parindent}{0pt}
\setlength{\parskip}{4pt}

\begin{document}

\begin{center}
  {\Large\bfseries Distributed Global Secondary Index Maintenance for Sharded PostgreSQL}\\[4pt]
  Kunj Bhesaniya (23B0995) \quad Dhruvraj Merchant (23B1041) \quad
  Dipen Sojitra (23B0913) \quad Harshil Solanki (23B1016)\\[2pt]
  {\small CS 631 \textperiodcentered{} Checkpoint 1 progress report \textperiodcentered{} 2 November 2026}
\end{center}

\section*{Problem}
When a table is sharded across PostgreSQL servers by primary key, a query on another column
(\texttt{WHERE email = ?}) must ask every shard, and \texttt{UNIQUE(email)} cannot be enforced because
each shard sees only its own rows. Spanner, CockroachDB and TiDB solve this with a \emph{global
secondary index} (GSI); core PostgreSQL and Citus have none. We build a GSI whose entry for key $k$
lives on node $\mathrm{hash}(k) \bmod N$, independent of where the row lives. Every write then touches
two nodes, so the core challenge is \textbf{keeping base table and index consistent under concurrent
writes and node crashes}.

\section*{Status}
The full synchronous GSI works end to end. We planned milestones M1--M6 for this checkpoint and
finished M1--M8.

\begin{center}
\begin{tabular}{lll}
\toprule
Part & Milestones & Status \\
\midrule
Setup: cluster, module, catalog          & M1, M2     & Done \\
Reads: plan selection, GsiLookup         & M8         & Done \\
Writes: interception, 2PC, uniqueness    & M3, M4, M5 & Done \\
Recovery: resolver, background worker    & M6, M7     & Done \\
Stress test and benchmarks               & M9         & Stress test done; benchmarks pending \\
Online index build, async mode           & Stretch    & After checkpoint \\
\bottomrule
\end{tabular}
\end{center}

\section{Setup: cluster, connections and GSI catalog}
All our code is one PostgreSQL C extension, \texttt{gsi\_fdw}: a modified copy of
\texttt{postgres\_fdw} (PostgreSQL's wrapper for remote servers) plus about 1,500 lines of our own code.
\begin{itemize}
  \item \textbf{Cluster:} Docker Compose, 1 coordinator and 3 data nodes (PostgreSQL 16). On the
        coordinator, \texttt{users} is hash-partitioned by \texttt{id}; each partition is a foreign
        table on one data node, so PostgreSQL's partition routing sends each row to its shard.
  \item \textbf{Connections:} we reuse postgres\_fdw's connection cache and per-node remote
        transactions, so index writes join the same transaction as the base write.
  \item \textbf{Catalog:} the coordinator stores the nodes (\texttt{gsi.gsi\_node}) and indexes
        (\texttt{gsi.gsi\_index}). \texttt{gsi\_create('users','email',unique)} creates
        \texttt{gsi\_users\_email(key, id, shard)} on every node. A table can have several GSIs,
        one per column.
  \item \textbf{Placement:} $\mathrm{hash}(k) \bmod N$ with the hash PostgreSQL uses for hash
        partitioning; one C function serves writes, reads and the checker.
  \item \textbf{One entry per row} \texttt{(key, id, shard)}: exact primary-key fetches, exact
        deletes, and a unique B-tree for global uniqueness.
\end{itemize}

\section{Reads: GSI plan selection and query routing}
A lookup on the indexed column asks one index node and only the shards that hold matching rows.
\begin{itemize}
  \item \textbf{Plan selection:} a planner hook (\texttt{set\_rel\_pathlist\_hook}) sees
        \texttt{WHERE email = <value>} (or any column with a ready GSI) and adds our plan, a
        \textbf{GsiLookup CustomScan}, next to PostgreSQL's ``scan all shards'' plan (Append).
  \item \textbf{Two hops:} (1) hash the key and ask that index node for \texttt{(id, shard)};
        (2) fetch \texttt{WHERE id = ANY(...)} from only those shards.
  \item \textbf{Recheck:} all original conditions are re-evaluated on fetched rows, so a stale
        index entry can never produce a wrong result.
  \item \textbf{Planner issue solved:} PostgreSQL rebuilds all plans of a partitioned table late in
        planning (\texttt{apply\_scanjoin\_target\_to\_paths}), which dropped our plan. After pruning
        we mark the table non-partitioned for that query, so our plan survives.
  \item \texttt{EXPLAIN} shows \texttt{Custom Scan (GsiLookup) on users} with the index node and
        number of shards fetched. \emph{Limit:} cost is a fixed heuristic for now.
\end{itemize}

\section{Writes: index maintenance and two-phase commit}
Every INSERT, UPDATE and DELETE changes the base row and its entry in every GSI of the table in one distributed
transaction, committed atomically with two-phase commit (2PC).
\begin{itemize}
  \item \textbf{Write interception:} for writes that can change an index entry (INSERT, DELETE,
        UPDATE of an indexed column) we disable postgres\_fdw's pushdown and batching, so our
        per-row hooks run; other UPDATEs keep the pushdown. UPDATE and DELETE also fetch the
        \emph{old} key of every indexed column, so we know which entries to remove.
  \item \textbf{2PC} (in postgres\_fdw's commit callback):
  \begin{enumerate}
    \item Before commit: \texttt{PREPARE TRANSACTION 'gsi\_<xid>\_<node>\_<role>'} on every node
          involved; any failure aborts everything.
    \item Decision: the coordinator's own local commit (its xid is in every prepared name).
    \item After commit: \texttt{COMMIT PREPARED} on \textbf{index nodes first, then base shards},
          so a visible row always has its index entry.
  \end{enumerate}
  \item \textbf{Global uniqueness:} a unique B-tree on the index node; a duplicate aborts the whole
        transaction, including entries already written to other GSIs. With two concurrent inserts of the same email, exactly one wins.
  \item \textbf{Not supported:} changing \texttt{id} (foreign partitions cannot move rows between
        shards, so we reject it instead of corrupting data).
\end{itemize}

\section{Recovery: finishing transactions after a crash}
A crash in the middle of 2PC leaves prepared transactions holding locks. Our resolver finishes them
by following the coordinator's recorded decision, as Citus (2PC recovery), CockroachDB (intent
resolution) and TiDB (lock resolution) do.
\begin{itemize}
  \item \textbf{Resolver:} reads \texttt{pg\_prepared\_xacts} on each node, asks the coordinator
        whether the transaction committed (\texttt{pg\_xact\_status}), then runs
        \texttt{COMMIT PREPARED} (index nodes first) or \texttt{ROLLBACK PREPARED}.
  \item \textbf{Background worker:} runs the resolver every 5\,s inside the coordinator; no manual step.
  \item \textbf{Fault injection:} a test-only setting \texttt{gsi.fail\_point} crashes the coordinator
        at an exact protocol step, since real crash windows last microseconds.
\end{itemize}

\begin{center}
\begin{tabular}{ll}
\toprule
Crash point & Result after recovery \\
\midrule
After PREPARE, before the decision     & Rolled back: row and entry absent \\
After the coordinator's commit         & Rolled forward: row and entry present \\
Between index commit and base commit   & Base finished: row present \\
Data node paused during commit         & Transaction aborts; resolver cleans up \\
\bottomrule
\end{tabular}
\end{center}

\section*{Limitations and plan after Checkpoint 1}
\textbf{Limitations:} no global snapshot (during the short commit window a reader may see the entry
before the row; the recheck keeps results correct); cross-node deadlocks are ended by a 5\,s lock
timeout; \texttt{gsi\_create} needs an empty table; one GSI per column (several per table); only \texttt{col = value}
uses the GSI (\texttt{IN}, ranges and \texttt{FOR UPDATE} ask every shard).

\textbf{Next steps:}
\begin{enumerate}
  \item \textbf{Benchmarks:} write throughput and latency with and without the GSI; lookup latency,
        GSI vs every shard, at 2, 3 and 4 shards.
  \item \textbf{Cost model:} selectivity-based cost, so the planner uses the GSI only when cheaper.
  \item \textbf{Online index build} on a populated table under live writes (F1 states: write-only,
        backfill, validate, ready).
  \item \textbf{Asynchronous mode} via logical decoding, with lag measurement; no global uniqueness
        in this mode.
\end{enumerate}

\end{document}
